\documentclass[11pt]{article}
\usepackage[margin=1in]{geometry}
\usepackage{amsmath,amssymb}
\usepackage{booktabs}
\usepackage{hyperref}

\title{Notes on panel dependence and alpha decay in ultra-high-dimensional drift localization}
\author{Working notes, FSDS benchmark}
\date{}

\begin{document}
\maketitle

\paragraph{Status of these notes.}
Statements are tagged as \emph{verified} (reproduced by code in this repository), \emph{heuristic} (derived or extrapolated, not tested), or \emph{planned} (not yet done).
The panel data-generating process (DGP) described in Section~4 is \textbf{not implemented yet}; all benchmark results in the repository use i.i.d.\ rows.

\section{Motivation}
In the quantitative-finance formulation of the benchmark, alpha decay is pure concept drift: after a regime change the distribution of the (cross-sectionally standardized) characteristics $X$ is unchanged, whereas $\mathbb{E}[Y\mid X]$ changes for a subset of characteristics.
Domain-classifier methods cannot see this (verified: random-forest domain classifier AUC $0.43$--$0.53$ in the pure concept-drift scan), so the relevant methods are those that look at the outcome regression (PermuCATE, LOCO-R-risk, SGShift).
The benchmark rows are i.i.d.\ with $n=500$ per regime.
Real panels have $N$ stocks observed over $T$ periods, and the rows are dependent.
These notes ask how that dependence changes the effective sample size, and hence the meaning of ``ultra-high dimensional'' ($p/n$) in this application.

\section{Sources of dependence}
Write $y_{it} = \sum_j b_{jt}\, g_j(x_{ijt}) + \lambda_i^\top g_t + e_{it}$ for stock $i$ and period $t$.
Table~\ref{tab:dep} lists the sources and their effect on the slope of a single characteristic.

\begin{table}[h]
\centering
\caption{Sources of dependence in a stock-month panel and their effect on drift detection.}
\label{tab:dep}
\begin{tabular}{p{3.2cm}p{4.6cm}p{6.2cm}}
\toprule
Source & Mechanism & Effect on slope estimation \\
\midrule
Common period shocks & Market-wide return $u_t$ & Negligible if characteristics are standardized or ranked within each period (absorbed by a period fixed effect) \\
Time-varying slopes & $b_{jt} = b_j + \tau\,\varepsilon_t$ shared by all stocks in period $t$ & Dominant: not averaged out by adding stocks (Section~3) \\
Cross-sectional residual correlation & Industry or size shocks $\lambda_i^\top g_t$ & Clustering by period, similar to time-varying slopes \\
Serial dependence & Persistent characteristics, overlapping return windows, persistent regimes & Matters only through autocorrelation of slopes or residuals (heuristic) \\
\bottomrule
\end{tabular}
\end{table}

\section{Effective sample size with time-varying slopes}
Consider a single standardized characteristic, $y_{it} = b_t x_{it} + e_{it}$ with $b_t = b + \tau\varepsilon_t$, $\operatorname{Var}(x)=1$ and $\operatorname{Var}(e)=\sigma^2$.
The pooled least-squares slope satisfies, approximately,
\begin{equation}
\operatorname{Var}(\hat b) \approx \frac{\sigma^2}{NT} + \frac{\tau^2}{T},
\qquad
\mathrm{deff} = 1 + \frac{N\tau^2}{\sigma^2},
\qquad
n_{\text{eff}} = \frac{NT}{\mathrm{deff}} \xrightarrow{N\to\infty} \frac{T\sigma^2}{\tau^2}.
\end{equation}
The second variance term does not depend on $N$.
Table~\ref{tab:mc} compares the formula with a Monte Carlo experiment (verified; \texttt{python/panel\_deff\_check.py}, output in \texttt{results/panel\_deff\_check.txt}; 800--1500 replicates, so the Monte Carlo error of the design effect is of the order of 5--10\%).

\begin{table}[h]
\centering
\caption{Design effect and effective sample size of the pooled slope ($b=0.3$, $\sigma=1$).}
\label{tab:mc}
\begin{tabular}{rrrrrrr}
\toprule
$N$ & $T$ & $\tau$ & $\phi_x$ & deff (MC) & deff (formula) & $n_{\text{eff}}$ \\
\midrule
500 & 12 & 0   & 0   & 1.0  & 1.0  & 6096 \\
500 & 12 & 0.1 & 0   & 6.2  & 6.0  & 968 \\
500 & 12 & 0.3 & 0   & 39.5 & 46.0 & 152 \\
2000 & 12 & 0.1 & 0  & 22.0 & 21.0 & 1093 \\
500 & 60 & 0.1 & 0   & 6.1  & 6.0  & 4907 \\
500 & 60 & 0.3 & 0   & 40.3 & 46.0 & 744 \\
500 & 12 & 0.1 & 0.9 & 5.7  & 6.0  & 1055 \\
500 & 60 & 0.1 & 0.9 & 5.9  & 6.0  & 5047 \\
\bottomrule
\end{tabular}
\end{table}

Three observations.
(i)~Quadrupling the number of stocks ($N=500\to2000$ at $T=12$) leaves $n_{\text{eff}}$ essentially unchanged ($968$ vs.\ $1093$); lengthening the sample does help ($T=12\to60$: $968\to4907$).
(ii)~Persistence of the characteristic ($\phi_x=0.9$) does not change the design effect when the slope shocks are independent over time.
(iii)~The formula is accurate for moderate $\tau$ and overstates the design effect by roughly 15\% at $\tau=0.3$.
If the slopes themselves are autocorrelated with coefficient $\phi_b$, the effective sample size is further divided by approximately $(1+\phi_b)/(1-\phi_b)$ (heuristic, not simulated).

\paragraph{Consequence for the benchmark (heuristic).}
The ratio $p/n$ in the benchmark should be read as $p/n_{\text{eff}}$.
A panel with $NT\sim10^5$ rows can have $n_{\text{eff}}$ of a few hundred to a few thousand, so the ultra-high-dimensional regime studied here remains the relevant one.
This was verified only for a single slope; the extension to sparse high-dimensional selection is an extrapolation.

\section{Proposed panel DGP (planned)}
\begin{enumerate}
\item \textbf{Characteristics.} Stock loadings are fixed over time; the factor component of each characteristic follows an AR(1) with $\phi_j\in[0.5,0.98]$; characteristics are rank-standardized within each period.
\item \textbf{Returns.} $y_{it} = \sum_j b_{jt} g_j(x_{ijt}) + \lambda_i^\top g_t + e_{it}$ with heavy-tailed $e_{it}$ (Student-$t_5$), optionally an MA$(h-1)$ component for overlapping horizons.
\item \textbf{Alpha decay.} For a random 5\% of the characteristics the slope changes at a known change point $T_0$: exponential decay with a given half-life, sign flip, or change of functional form. The regime label is $W_t=\mathbf 1\{t>T_0\}$.
\item \textbf{Free parameters.} $\tau$ (slope volatility), $\phi_b$ (slope persistence), $T_{\text{old}}$, $T_{\text{new}}$ and $N$, calibrated to the dispersion of cross-sectional (Fama--MacBeth) slope estimates rather than chosen arbitrarily.
\end{enumerate}

\section{Implications for the methods}
\begin{itemize}
\item \textbf{Cross-fitting must split by time blocks with an embargo, not by rows.} Rows in the same period share $g_t$ and adjacent periods share persistent characteristics, so a random row split leaks information and inflates nuisance $R^2$ and selection AUC. The i.i.d.\ results in the repository are optimistic in this respect (heuristic).
\item \textbf{Permutations must be performed within periods}, with the conditional sampler also fitted within periods, to preserve the period structure.
\item \textbf{Detection power is governed by $\tau$ and $T$, not by $N$.} For a single characteristic the standard error of the before--after slope difference is of order $\sqrt{\tau^2(1/T_{\text{old}}+1/T_{\text{new}})}$ plus an $N$-dependent term. With $\tau=0.1$ and $T_{\text{old}}=T_{\text{new}}=36$ this is about $0.024$, so a change of the order of $0.12$ (five standard errors, as a rough multiplicity allowance for $10^4$ characteristics) would be detectable (heuristic arithmetic, not simulated).
\end{itemize}

\section{Limitations}
\begin{enumerate}
\item The panel DGP is not implemented. Everything above about sparse high-dimensional selection under dependence is a hypothesis; the Monte Carlo evidence covers one slope.
\item The benchmark results use i.i.d.\ rows, $n=500$ per regime, a single data-generating process, and Student-$t_5$ noise without stochastic volatility.
\item The main pure concept-drift run has five grid points with one replicate each; the $p/n$ scan has three replicates (standard deviations $0.1$--$0.15$ at small $p/n$); the PermuCATE nuisance study has five replicates (standard errors $0.02$--$0.06$ in AUC). Differences smaller than about $0.05$--$0.1$ AUC cannot be resolved.
\item The regime change point is known. In practice it must be estimated, and an error in $T_0$ blurs the contrast.
\item In the nuisance study, the explanation that limited outcome-regression accuracy drives the loss is supported by ablations, but the oracle also degrades with $p/n$, so part of the decline is due to the CATE and permutation stages. The noisy-oracle ablation uses independent Gaussian errors, which differ from the structured errors of real estimators.
\item No real data, transaction costs, look-ahead control or turnover constraints are considered. The calibration of $\tau$ and of the factor structure to real data has not been done.
\item The effective-sample-size formula assumes a single slope, independent slope shocks over time and homoskedastic noise.
\end{enumerate}

\section{Planned improvements}
\begin{enumerate}
\item \textbf{Implement the panel DGP and block cross-fitting} (Sections~4--5) and quantify the optimism of the row-wise split by comparing the two on the same data.
\item \textbf{Scan $p/n_{\text{eff}}$} over $(N,T,\tau,\phi_b)$ and test the prediction that, at fixed $p$, the AUC depends mainly on $T$ and $\tau$ and only weakly on $N$.
\item \textbf{Period-level inference.} Estimate period-by-period cross-sectional drift statistics and test them across periods (Fama--MacBeth style, with HAC standard errors), which uses the correct unit of replication.
\item \textbf{Calibrated error control.} Obtain null distributions for the selection scores with block permutations of the regime label (circular shifts of $T_0$) or a period-level block bootstrap, and report false-discovery proportions rather than AUC alone.
\item \textbf{Hierarchical or grouped detection.} Detect drift at the level of characteristic families first and then within families, using sample splitting to avoid selection bias, as per-characteristic attribution is statistically unrealistic at high $p/n_{\text{eff}}$.
\item \textbf{Unknown change point.} Add a scan over candidate $T_0$ (CUSUM-type statistic) and report the sensitivity of the selection to its error.
\item \textbf{Better replication.} Use more replicates and several seeds for the main pure concept-drift run, with confidence intervals for the AUC.
\item \textbf{Nuisance models.} Pool the two arms in a single outcome model with the regime indicator as a covariate, increase the number of cross-fitting folds, and neutralize the factor component. A first factor-augmented variant gave a mixed result (an AUC gain only at $p/n=100$) and is treated as exploratory.
\item \textbf{Real-data check.} Apply the pipeline to a public set of anomaly characteristics with documented post-publication decay and check whether the flagged characteristics agree with the documented decay.
\end{enumerate}

\end{document}
